\section{The Engine}\label{sec:approach}

The engine is a level-synchronous breadth-first search over the zone graph in
which every thread inserts into one shared passed set.  We describe it bottom-up:
first the passed set and the operation a thread performs on it, then the parallel
level loop that calls that operation and the protocol under which it is called,
then the work that runs at a level barrier.  Correctness and cost bounds are
deferred to \cref{sec:correct}; this section states what the engine does and why
each mechanism is there.

\subsection{The passed set and its update}\label{sec:passed}

The passed set is one structure that every thread may touch, partitioned into $S$
\emph{stripes} by a hash of the discrete state.  A stripe owns a lock, a slot
table mapping a discrete state to the head of a chain, and the chain nodes
themselves; a node holds the index of one stored zone and the breadth-first level
at which it was stored.  Since $\stripe$ is a function of the discrete state, all
updates at one discrete state take one lock, and updates at discrete states in
different stripes touch disjoint memory.

Zones are not held in the chains.  All distinct zones live in one \emph{zone
table} $\zt$ shared by every stripe, which stores each zone once as a packed
record and hands out an index; chain nodes and frontier entries hold indices.
Interned records are immutable while referenced, because every operation below
mutates only the incoming zone and reads stored ones, so reads take no lock and
only the first sight of a distinct zone takes one.

\begin{algorithm}[t]
\caption{The stored-set update and the coverage probe at one discrete state}
\label{alg:update}
\begin{algorithmic}[1]
\Procedure{Update}{$q,\Zone,k$} \Comment{exclusive lock on $\stripe(q)$ held}
  \If{$\exists\, Y \in \Passed(q):\ \Zone \subseteq \mg(Y)$}
      \Return \textsc{dropped} \Comment{on the packed records}
  \EndIf
  \State $A \leftarrow \Passed(q)$ as dense matrices, newest first
         \Comment{mirror, else one unpack each}
  \If{$\exists\, Y \in A:\ \Zone \subq_q Y$} \Return \textsc{dropped}
      \Comment{(a)} \EndIf
  \State retire $\{Y \in A : Y \subq_q \Zone\}$ and remove it from $A$
         \Comment{(b)}
  \While{some $Y$ among the $m$ newest of $A$ merges exactly with $\Zone$}
    \State $\Zone \leftarrow \hull(\Zone,Y)$; retire $Y$ and remove it from $A$
           \Comment{(c)}
  \EndWhile
  \State intern $\Zone$, add it to $\Passed(q)$ stamped $k$, maintain the mirror
  \State \Return \textsc{inserted}
\EndProcedure
\Statex
\Procedure{Covered}{$q,\Zone$} \Comment{shared lock on $\stripe(q)$ held; writes nothing}
  \If{$\Passed(q)$ has a mirror}
      \Return $\exists\, Y \in \Passed(q):\ \Zone \subq_q Y$ \EndIf
  \State \Return $\exists\, Y \in \Passed(q):\ \Zone \subseteq \mg(Y)$
\EndProcedure
\end{algorithmic}
\end{algorithm}

\Cref{alg:update} states the two procedures a thread calls per successor, where
$\Passed(q)$ ranges over live zones only.  \textsc{Update} is the operation the
search needs: fold~(a) drops the candidate $\Zone$ when a stored zone subsumes
it, fold~(b) retires the stored zones that $\Zone$ subsumes, and fold~(c) absorbs
stored zones that merge exactly with $\Zone$, replacing $\Zone$ by the hull and
repeating, because a hull can merge with a zone the original could not.  What
survives is interned and added.  \textsc{Covered} is the read-only prefix of that
operation: it decides fold~(a) and writes nothing.

Three parts of \cref{alg:update} exist to keep the cost of one call low, and each
is stated here and justified in \cref{sec:cost}.

\emph{The packed-record filter, line~2.}  Every predicate in folds (a)--(c) reads
a stored zone as a dense matrix, and stored zones are packed records, so a direct
implementation unpacks a record per probe at cost $O(n^3)$.  A minimal constraint
system is a generating set, and inclusion can be decided against generators
without rebuilding what they generate: $\Zone \subseteq \mathrm{close}(\mg)$
holds exactly when $\Zone[i][j] \le b$ for every $(i,j,b) \in \mg$.  Line~2 runs
that test, which costs $O(|\mg|)$ and needs no scratch matrix.  A hit ends the
call; a miss falls through to the unchanged folds.  \textsc{Covered} uses the
same test, and only that test, when the stored set has no mirror.

\emph{The dense mirror.}  Where line~2 does not fire, the folds read every live
zone densely, so an update costs $\Theta(L \cdot n^3)$ for a stored set of $L$
live zones.  Stored sets of size close to one, the common case on process
networks, pay little; stored sets that grow do not, and a model of one automaton
with one location and $N$ independent clocks keeps its entire passed set in a
single stored set of thousands of mutually incomparable zones.  Such a stored set
additionally keeps its live zones as dense canonical matrices in insertion order,
in a \emph{dense mirror} that the folds read in place.  The packed records remain
the storage of record; the mirror is a cache, appended on insertion, its slot
cleared on retirement, rebuilt when half its slots are cleared, and subject to a
byte budget that may refuse to create it.  It is created once the unpacking the
stored set has already cost crosses a threshold: with $v$ the number of updates
that have walked it, when $v \cdot L \ge W$ and $L \ge L_{\min}$.  What a mirror
saves is one unpack per live zone \emph{per visit}, so the product governs it and
a rule on $L$ alone cannot see the difference between $16$ zones walked $100$
times and $64$ walked once.  The floor $L_{\min}$ is a memory condition: a mirror
of a one-zone stored set is a private copy of a zone the zone table was sharing.

\emph{The merge window $m$.}  Fold~(c) restarts after each success, so its cost is
quadratic in $L$ --- and it is quadratic exactly where stored sets are large.  A
candidate therefore attempts to merge only with the $m$ most recently inserted
live zones.  Zones that merge are almost always recent neighbours, because a
merge happens when a successor extends a zone stored moments earlier.

\subsection{Parallel level-synchronous exploration}\label{sec:levels}

\figengine

A search that interleaves a call to $\succfn$ with the update of its result
serialises, because the update performed for one successor is observed by the
check performed for the next.  Grouping the search by breadth-first level breaks
that dependency for the calls to $\succfn$.  The \emph{frontier} $\Frontier_k$ is
the set of symbolic states inserted during level $k$; every one of them was
checked against the passed set before being placed there, so their expansions are
independent and may run in any order.  A frontier entry is a (discrete state,
zone index) pair: its zone is read from $\zt$ when the entry is expanded, rather
than copied into the frontier, so the frontier does not hold a second copy of
zones the passed set already stores.

\Cref{alg:engine} states one level and \cref{fig:engine} shows it.  Threads claim
chunks of the frontier under a dynamic schedule and insert each successor into
the stripe of its discrete state as soon as it is generated.  Fusing expansion
and insertion removes a second pass over data that has just left cache, and it
lets every thread see, within a level, the zones the others have stored.

\begin{algorithm}[t]
\caption{One level of the parallel engine}
\label{alg:engine}
\begin{algorithmic}[1]
\Require network $\net$, predicate $\varphi$, thread count $T$, stripe count $S \gg T$
\State $\Passed(q_0) \leftarrow \{\Zone_0\}$;\enspace
       $\Frontier \leftarrow [(q_0,\Zone_0)]$;\enspace $k \leftarrow 0$
\While{$\Frontier \neq \emptyset$}
  \State $k \leftarrow k+1$
  \ParFor{$i = 0,\dots,|\Frontier|-1$} \Comment{dynamic schedule}
    \State $(q,\Zone) \leftarrow$ decode $\Frontier[i]$
    \ForAll{$(q',\Zone') \in \succfn(q,\Zone)$}
      \If{$\varphi(q',\Zone')$} \Return \textbf{true} \EndIf
      \State \textbf{lock} $\stripe(q')$ \textbf{shared};\enspace
             $c \leftarrow \Covered(q',\Zone')$;\enspace \textbf{unlock}
      \If{$c$} \textbf{continue} \EndIf
      \State \textbf{lock} $\stripe(q')$ \textbf{exclusive};\enspace
             $r \leftarrow \Add(q',\Zone',k)$;\enspace \textbf{unlock}
      \If{$r = \textsc{inserted}$} record $q'$ in this thread's touched set \EndIf
    \EndFor
  \EndParFor
  \State $\Frontier \leftarrow$ zones stamped $k$ of the touched discrete states
  \State compact retired nodes; reclaim unreferenced records of $\zt$; revise the abstraction
\EndWhile
\Return \textbf{false}
\end{algorithmic}
\end{algorithm}

Two details of \cref{alg:engine} are load-bearing for \cref{sec:correct}.  The
predicate $\varphi$ is evaluated on \emph{generated} symbolic states, on line~7,
and never on a stored zone, so a zone that fold~(c) built as a convex union is
never offered to it.  And a thread records the discrete states it inserted into
rather than the states themselves; the next frontier is read back out of the
passed set at the barrier, so a zone inserted early in a level and retired later
in the same level does not enter it.

\paragraph{The two-phase stripe protocol.}
Lines~8--11 split the operation of \cref{alg:update} into its read-only and its
mutating half.  \textsc{Covered} runs under a \emph{shared} stripe lock, so
probes issued by different threads on one stripe proceed concurrently, and only a
thread whose probe answers \emph{no} takes the \emph{exclusive} lock and runs
\textsc{Update} (\cref{fig:protocol}).  The split pays because the operation is
read-mostly and because collisions are structural rather than random: a successor
is routed to the stripe of its discrete state, and a symmetric network sends many
successors of one level to the same discrete state, hence to one stripe.  It is
sound because a positive probe cannot be invalidated while the shared lock is
released, which is \cref{lem:monotone}.  A probe answering \emph{no} where
\textsc{Update} would have dropped the candidate costs one exclusive lock and
nothing else, since \textsc{Update} repeats the test.

\figprotocol

Critical sections here last microseconds, so each stripe carries an eight-byte
reader/writer lock that spins a bounded number of times before yielding, with the
bound set from the thread count relative to the core count.

\paragraph{The interface to the correctness argument.}
\Cref{sec:correct} uses the update only through the following contract.  Write
$\Passed(q)$ for the stored set before an update with candidate $\Zone$ at $q$,
and $\Passed'(q)$ for the stored set after it.
\emph{(U1)} If $\mathrm{cov}_\Passed(q,\Zone)$ then $\Passed'(q) = \Passed(q)$
and nothing is recorded as inserted.
\emph{(U2)} Otherwise $\Passed'(q) = \{\Zone^{*}\} \cup (\Passed(q) \setminus R)$
for some $R \subseteq \Passed(q)$ and some $\Zone^{*}$ with $\Zone \subseteq
\Zone^{*} \subseteq \Zone \cup \bigcup R$ and $Y \subq_q \Zone^{*}$ for every $Y
\in R$.
\emph{(U3)} All updates at one discrete state are totally ordered, and each
observes the stored set left by its predecessor.

\subsection{The level barrier}\label{sec:barrier}

Expansion is parallel and contended only where discrete states collide, so
whatever runs between levels is the engine's serial term, and it touches the
passed set, which is larger at every level.  Four things run there.

\emph{Collecting the next frontier.}  The next frontier is the set of zones
stamped with the level just finished, read out of the passed set over the
discrete states the threads recorded.  A discrete state may be recorded by
several threads and must be scanned once, so the records are deduplicated: the
keys are partitioned by a hash into $B$ buckets and each bucket is deduplicated
and scanned independently, in one parallel pass.  The bucket function is taken
from the same hash as $\stripe$, so a bucket owns a contiguous range of stripes
and its scan stays within one slice of the passed set.

\emph{Compacting retired nodes.}  Folds (b) and (c) retire chain nodes by marking
them, because unlinking while other threads hold chain positions is not safe
during a level.  Marked nodes are traversed and skipped by every later chain walk,
so they are removed between levels, one thread per stripe, by rebuilding each
chain from its live nodes.  A stripe is rebuilt once its marked nodes reach a
third of its nodes, and that count is maintained incrementally by the two sites
that retire a node rather than recomputed by a scan, so the trigger is $O(1)$ and
can be evaluated at every barrier.

\emph{Reclaiming the zone table.}  A retired node may hold the last reference to
an interned record.  Unreferenced records are reclaimed by marking every record
named by a live chain node or by a frontier entry --- the frontier is a root ---
and returning the rest to a free list.  The trigger is growth-based.  Let $A$ be
the number of records currently handed out, that is the high-water index count
minus the size of the free list; $A$ rises with every fresh intern and falls back
to the live set after a reclamation, so it is the quantity a trigger must watch.
A reclamation runs when $A \ge \mathit{next}$, after which $\mathit{next}
\leftarrow \max(2A_{\mathrm{live}}, \tau)$ for a floor $\tau$ below which the
scan is not worth its own cost.  Freed records and freed indices are recycled in
place, so no index held by a chain node is ever renumbered; the price is that
every memo mapping record content to an index, or an index to a decoded matrix,
must be invalidated, which the engine does by stamping each memo with a
reclamation counter.

\emph{Revising the abstraction.}  Subsumption and extrapolation are two routes to
finiteness and they act differently: $\subq_q$ only prunes, dropping a zone when
a stored one simulates it, whereas $\extraLU$ widens, mapping many zones onto one
representative.  Where the reachable zones at a discrete state are mutually
$\subq_q$-incomparable, pruning achieves nothing and widening collapses them;
where they collapse under $\subq_q$ anyway, widening only costs precision.  Which
case holds is a property of the model, and the distinguishing signal is available
in $O(1)$ per stripe at a barrier: the ratio of live stored zones to distinct
discrete states, near one when subsumption is collapsing the stored sets and
climbing when they fill with incomparable zones.  The engine therefore starts
without extrapolation and switches it on, for every zone generated thereafter, at
the first barrier at which the ratio crosses a threshold $\theta$.  The switch is
monotone, so a run performs at most one.
